\section{Dataset and Preprocessing}
\label{sec:dataset}

\subsection{Dataset Description}

This study utilizes the Oral Diseases dataset from Kaggle~\cite{kaggle_oral_diseases} to evaluate our framework for automated oral disease classification. The dataset comprises \textbf{12,817} intraoral images categorized into six common oral pathologies.

\textbf{Dataset Source}: The Oral Diseases dataset\footnote{\url{https://www.kaggle.com/datasets/salmansajid05/oral-diseases}} contains clinical photographs of various oral conditions collected from dental examinations. The dataset covers six disease categories: Calculus, Caries, Gingivitis, Hypodontia, Tooth Discoloration, and Ulcers. Each image is labeled with a single disease category, enabling supervised learning for multi-class classification tasks.

\textbf{Supplementary Data Sources}: In addition to the primary Oral Diseases dataset, we reference the following publicly available datasets for potential future expansion and comparative research: the Teeth Dataset\footnote{\url{https://www.kaggle.com/datasets/rajapriyanshu/teeth-dataset}} containing dental images for teeth classification tasks, and the Oral Infection dataset\footnote{\url{https://www.kaggle.com/datasets/sizlingdhairya1/oral-infection}} providing labeled images of oral infection conditions. These datasets enable cross-dataset validation and domain adaptation studies in oral disease recognition.

\textbf{Dataset Composition}: The complete dataset of 12,817 images is divided into training, validation, and test sets with the following distribution across six disease categories: Calculus (1,296 images, 10.11\%), Caries (2,382 images, 18.58\%), Gingivitis (3,358 images, 26.20\%), Hypodontia (1,406 images, 10.97\%), Tooth Discoloration (1,834 images, 14.31\%), and Ulcers (2,541 images, 19.83\%). The dataset exhibits natural class imbalance reflecting the relative prevalence of these conditions in clinical practice.

\begin{table}[htbp]
\centering
\caption{Oral Diseases Dataset Statistics}
\label{tab:dataset_stats}
\resizebox{0.48\textwidth}{!}{%
\begin{tabular}{lcccc}
\toprule
\textbf{Disease Category} & \textbf{Train} & \textbf{Val} & \textbf{Test} & \textbf{Total} \\
\midrule
Calculus & 932 (10.44\%) & 234 (9.33\%) & 130 (9.38\%) & 1,296 (10.11\%) \\
Caries & 1,714 (19.21\%) & 429 (17.11\%) & 239 (17.24\%) & 2,382 (18.58\%) \\
Gingivitis & 2,161 (24.22\%) & 778 (31.03\%) & 419 (30.23\%) & 3,358 (26.20\%) \\
Hypodontia & 969 (10.86\%) & 278 (11.09\%) & 159 (11.47\%) & 1,406 (10.97\%) \\
Tooth Discoloration & 1,320 (14.79\%) & 330 (13.16\%) & 184 (13.28\%) & 1,834 (14.31\%) \\
Ulcers & 1,828 (20.48\%) & 458 (18.27\%) & 255 (18.40\%) & 2,541 (19.83\%) \\
\midrule
\textbf{Total} & \textbf{8,924} & \textbf{2,507} & \textbf{1,386} & \textbf{12,817} \\
\bottomrule
\end{tabular}
}
\end{table}

\subsection{Disease Categories}

\textbf{Calculus}: Dental calculus (tartar) is mineralized dental plaque that forms on tooth surfaces and can lead to periodontal disease. The dataset contains 1,296 calculus images showing various degrees of calculus buildup on teeth.

\textbf{Caries}: Dental caries (tooth decay) represents bacterial infection causing demineralization and destruction of tooth structure. With 2,382 images, this is the second most prevalent condition in the dataset.

\textbf{Gingivitis}: Inflammation of the gingival tissue caused by bacterial plaque accumulation. As the most represented condition with 3,358 images, gingivitis manifests as redness, swelling, and bleeding of the gums.

\textbf{Hypodontia}: Congenital absence of one or more teeth, excluding third molars. The dataset includes 1,406 images documenting various patterns of missing teeth.

\textbf{Tooth Discoloration}: Abnormal tooth coloration resulting from intrinsic or extrinsic factors including medications, trauma, or dietary habits. This category comprises 1,834 images.

\textbf{Ulcers}: Oral ulcerations appearing as painful lesions on the oral mucosa. The dataset contains 2,541 images showing various types and locations of oral ulcers.

\subsection{Annotation Protocol and Limitations}

\textbf{Annotation Process}: Each image in the dataset is labeled with a single disease category. According to the dataset documentation, annotations were performed by dental professionals based on visual examination of intraoral photographs. The labeling follows a single-label paradigm where each image is assigned to exactly one of six disease categories.

\textbf{Limitations and Considerations}: As a community-contributed Kaggle dataset, the Oral Diseases dataset has several limitations that should be acknowledged:

\begin{itemize}
    \item \textbf{No reported Inter-Annotator Agreement (IAA)}: Unlike peer-reviewed clinical datasets such as DENTEX~\cite{hamamci2023dentex}, the dataset does not provide inter-rater reliability metrics (e.g., Cohen's Kappa, Fleiss' Kappa) to quantify annotation consistency.
    \item \textbf{Single annotator per image}: The annotation protocol does not specify whether multiple annotators independently labeled each image, which is standard practice in clinical benchmark datasets.
    \item \textbf{Limited demographic information}: Patient demographics, imaging equipment specifications, and clinical settings are not documented.
    \item \textbf{No external validation}: The dataset has not undergone peer-reviewed validation or independent quality assessment.
\end{itemize}

\textbf{Quality Control Measures}: To mitigate these limitations, we performed manual inspection of a random sample (5\%) of images from each class to verify label correctness. Additionally, we excluded visibly corrupted or unrelated images during preprocessing. Despite these measures, annotation noise may still affect model performance and should be considered when interpreting results.

\subsection{Comparison with Benchmark Datasets}

Table~\ref{tab:dataset_comparison} compares the Oral Diseases dataset with established dental imaging benchmarks. While the dataset provides a reasonable sample size and class diversity, it lacks the rigorous annotation protocols of peer-reviewed datasets.

\begin{table}[htbp]
\centering
\caption{Comparison with established dental imaging datasets}
\label{tab:dataset_comparison}
\resizebox{0.48\textwidth}{!}{%
\begin{tabular}{lccccc}
\toprule
\textbf{Dataset} & \textbf{Size} & \textbf{Classes} & \textbf{Modality} & \textbf{IAA} & \textbf{Peer-Rev.} \\
\midrule
DENTEX~\cite{hamamci2023dentex} & 3,163 & 4 & Panoramic & $\kappa$=0.85 & \checkmark \\
Pediatric Panoramic~\cite{alsakar2024multilabel} & 2,500 & 5 & Panoramic & $\kappa$=0.82 & \checkmark \\
QLF Caries~\cite{lee2017qlf} & 1,500 & 2 & QLF & Reported & \checkmark \\
\textbf{Oral Diseases (Ours)} & \textbf{12,817} & \textbf{6} & \textbf{Intraoral} & \textbf{N/A} & $\times$ \\
\bottomrule
\end{tabular}
}
\end{table}

\textbf{Justification for Dataset Selection}: Despite these limitations, we selected the Oral Diseases dataset for several reasons: (1) it is the largest publicly available intraoral photograph dataset for oral disease classification, (2) it covers six clinically relevant disease categories with natural class imbalance, and (3) it enables reproducibility and comparison by other researchers. We acknowledge these limitations and recommend future work on curated clinical datasets with standardized annotation protocols.

\subsection{Image Preprocessing and Augmentation}

All oral disease images undergo standardized preprocessing to normalize characteristics and enhance model performance:

\begin{enumerate}
    \item \textbf{Image Loading}: Load RGB images and maintain original aspect ratios during initial processing
    \item \textbf{Resizing}: Standardize all images to $224\times224$ pixels for compatibility with pre-trained CNN architectures
    \item \textbf{Intensity Normalization}: Rescale pixel values from [0, 255] to [0, 1] range
    \item \textbf{Standardization}: Apply Z-score normalization using ImageNet statistics (mean=[0.485, 0.456, 0.406], std=[0.229, 0.224, 0.225]) to leverage transfer learning
\end{enumerate}

To enhance model robustness and generalization across clinical variations, we apply comprehensive data augmentation during training:

\begin{itemize}
    \item \textbf{Geometric}: Random rotation ($\pm 15^{\circ}$), scaling (0.8--1.2$\times$), translation ($\pm 10\%$), horizontal flipping
    \item \textbf{Photometric}: Random brightness/contrast adjustment ($\pm 20\%$), hue/saturation shift ($\pm 10\%$), Gaussian noise ($\sigma=0.01$)
    \item \textbf{Advanced}: Random cropping with padding, cutout (random erasing with $16\times16$ patches), and color jittering
\end{itemize}

These augmentations simulate natural variations in clinical photography including different lighting conditions, camera angles, and image quality, improving the model's ability to generalize to real-world deployment scenarios.

\subsection{Data Splitting Strategy}

The dataset is pre-divided into training (69.63\%), validation (19.56\%), and test (10.81\%) sets with 8,924, 2,507, and 1,386 images respectively. This splitting strategy ensures adequate data for model training while maintaining sufficient samples for validation and final evaluation.

The training set is used for model parameter optimization, the validation set guides hyperparameter tuning and early stopping decisions, and the test set remains held-out until final evaluation to provide unbiased performance estimates. All results reported in this paper use test set performance to ensure fair comparison.

The dataset exhibits natural class imbalance, with Gingivitis being the most prevalent (26.20\%) and Calculus the least common (10.11\%). This imbalance reflects real-world clinical distributions but requires careful consideration during model training through techniques such as class weighting, focal loss, or resampling strategies.

\begin{figure}[htbp]
\centering
\fbox{\parbox{0.9\columnwidth}{\centering\vspace{2cm}[Sample images from oral diseases dataset will be inserted here]\vspace{2cm}}}
\caption{Sample intraoral images from the Oral Diseases dataset showing six disease. Images demonstrate the visual variability within each disease category.}
\label{fig:dataset_samples}
\end{figure}